% SPDX-FileCopyrightText: Copyright (c) 2026 NVIDIA CORPORATION & AFFILIATES. All rights reserved.
% SPDX-License-Identifier: MIT
% Compile with LuaLaTeX. Content matches the ReportLab-rendered technical_report.pdf.
\documentclass[10pt]{article}
\usepackage[margin=0.75in]{geometry}
\usepackage{fontspec,longtable,array,booktabs,xcolor,hyperref}
\setmainfont{TeX Gyre Heros}
\definecolor{reportblue}{HTML}{17627C}
\hypersetup{colorlinks=true,urlcolor=reportblue}
\setlength{\parindent}{0pt}
\setlength{\parskip}{0.5em}
\setlength{\emergencystretch}{3em}
\begin{document}

TECHNICAL REPORT\par

{\LARGE\bfseries Combining Label Deconvolution and Frozen Walk-Text Embeddings for Product Classification\par}

A validation-locked, ten-seed study with SAGN/SLE and Correct \& Smooth\par

\textbf{Meghana Shanthappa · Junhao Shen · Rishi Puri}\\ NVIDIA\par

{\small Original report: September 8, 2026 · Editorial revision: November 2, 2026\par}

\subsection*{Abstract}

We study node classification on \textbf{ogbn-products} by combining released Label Deconvolution (LD) BERT features, frozen EmbeddingGemma encodings of graph-walk text, and original OGB features. The resulting 1,636-dimensional node representation feeds a four-stage SAGN model with self-label-enhanced training (SLE), followed by Correct \& Smooth (C\&S). Across ten fresh seeds, the validation-locked v2 pipeline achieves \textbf{94.0584\% ± 0.0433\% validation accuracy} and \textbf{87.4984\% ± 0.0490\% test accuracy} (mean ± sample standard deviation).\par

\section*{1. Dataset and study design}

We use the official ogbn-products split and the OGB 1.3.6 evaluator. The dataset contains 2,449,029 nodes, 47 classes, and 61,859,140 undirected product relations, represented after loading as 123,718,280 directional edge records.\par

\begin{longtable}{@{}p{2.213in}p{2.213in}p{2.213in}@{}}

\toprule

Training nodes & Validation nodes & Test nodes \\

\midrule

196,615 & 39,323 & 2,213,091 \\

\bottomrule\end{longtable}

V2 uses seeds 10-19 and reports all ten runs without selection. Seeds 0-9 were used during v1 exploration. Sample standard deviations use the n - 1 denominator. \textbf{External data: Yes.} [1, 7]\par

\begin{longtable}{@{}p{6.88in}@{}}

\toprule

\textbf{Scope of the result.} The reported accuracy belongs to the complete LD + frozen walk-text + raw features + SAGN/SLE + C\&S pipeline. \\

\bottomrule\end{longtable}

\newpage

\section*{2. Method}

\subsection*{2.1 Node representation}

\begin{longtable}{@{}p{2.033in}p{0.783in}p{3.824in}@{}}

\toprule

Component & Dim. & Source \\

\midrule

LD BERT features & 768 & Released products\_sagn/hidden\_state.pt \\

Frozen walk-text features & 768 & Mean-pooled, L2-normalized EmbeddingGemma encodings \\

Original OGB features & 100 & Original ogbn-products node features \\

Concatenated input & 1,636 & Order: LD, walk-text, raw \\

\bottomrule\end{longtable}

LD features are reused from the released feature tensor; LD is not retrained in this experiment. The concatenated features are then used to train the downstream SAGN/SLE predictor. [2, 3]\par

\subsection*{2.2 WalkLM and Label Deconvolution: a product example}

\textbf{Example product graph.} In an illustrative co-purchase graph, headphones connect to earbuds and a phone case; the earbuds connect to Sony WF-1000XM5 earbuds. Products are nodes, descriptions are attributes, and links indicate co-purchases.\par

\textbf{WalkLM: read product text along a graph walk.} A graph walk follows connected nodes, for example \textbf{Headphones → Earbuds → Sony WF-1000XM5 earbuds}. WalkLM converts the visited products' attributes into text and fine-tunes a language model on these sequences. Its embeddings combine product meaning with graph context. Our implementation uses walk-text with a frozen encoder. [4]\par

\textbf{LD (Label Deconvolution): prepare text features for graph aggregation.} A GNN mixes the headphones' features with neighboring audio products and a phone case, although connected products may have different categories. LD uses training-label information propagated over the graph to learn adjusted training targets for the text encoder, helping it produce features suited to that later mixing. The headphones' correct category remains unchanged. We reuse released LD-trained BERT embeddings. [3]\par

The encoder is \textbf{google/embeddinggemma-300m}, loaded with AutoModel. Its final hidden states are mean-pooled over non-padding tokens using the attention mask, then L2-normalized. The encoder is frozen. Unlike original WalkLM, this experiment does not fine-tune a language model on graph-walk text; \textbf{WalkLM-style} refers only to graph-walk textualization. [4, 6]\par

\subsection*{2.3 Predictor and post-processing}

\begin{longtable}{@{}p{3.491in}p{3.269in}@{}}

\toprule

SAGN/SLE configuration & Value \\

\midrule

Feature hops / label hops & 2 / 9 \\

Hidden channels / attention heads & 512 / 1 \\

Learning rate / batch size & 0.001 / 50,000 \\

Dropout values (reported order) & 0.5 / 0.4 / 0.2 / 0.5 \\

Epochs across four stages & 300 / 100 / 100 / 100 \\

Self-training thresholds & 0.95 / 0.95 / 0.999 \\

\bottomrule\end{longtable}

{\small C\&S uses 50 correction iterations and 50 smoothing iterations, correction alpha 0.35, smoothing alpha 0.60, and scale 20. Only training labels enter SAGN/SLE label features and C\&S propagation. Validation labels support early stopping and selection; test labels are excluded from propagation. [2, 5, 7]\par}

\newpage

\section*{3. Results and interpretation}

All accuracies below are percentages, reported as mean ± sample standard deviation. V1 is exploratory and test-aware; v2 uses the later validation-locked protocol. [7]\par

\begin{longtable}{@{}p{3.074in}p{1.783in}p{1.783in}@{}}

\toprule

Evaluation & Validation (\%) & Test (\%) \\

\midrule

V1 + C\&S: exploratory & 94.0010 ± 0.0299 & 87.4753 ± 0.0487 \\

V2 final stage before C\&S & 93.9204 ± 0.0522 & Not evaluated \\

V2 validation-selected C\&S & \textbf{94.0584 ± 0.0433} & \textbf{87.4984 ± 0.0490} \\

\bottomrule\end{longtable}

The source code indexes the four training stages as 0-3; the pre-C\&S row refers to stage 3, the fourth and final stage. Test accuracy was not computed for that row.\par

\subsection*{Leaderboard comparison}

The official ogbn-products leaderboard, checked on November 2, 2026, lists \textbf{GLEM+GIANT+SAGN+SCR} first with \textbf{87.37\% ± 0.06\% test accuracy}. The table below compares its leading entries with our reported result. All entries use external data.\par

\begin{longtable}{@{}p{3.491in}p{1.574in}p{1.574in}@{}}

\toprule

Method / status & Test accuracy (\%) & Validation (\%) \\

\midrule

Our complete pipeline (reported) & \textbf{87.4984 ± 0.0490} & 94.0584 ± 0.0433 \\

Published \#1: GLEM+GIANT+SAGN+SCR & 87.37 ± 0.06 & 94.00 ± 0.03 \\

Published \#2: LD+GIANT+SAGN+SCR & 87.18 ± 0.04 & 93.99 ± 0.02 \\

Published \#3: GraDBERT+GIANT \& SAGN+SLE+CnS & 86.92 ± 0.07 & 93.71 ± 0.03 \\

\bottomrule\end{longtable}

Our mean test accuracy exceeds the published \#1 score by \textbf{0.1284 percentage points} (87.4984\% - 87.37\%). Our pipeline combines LD-trained product-text features, frozen walk-text features, and raw OGB features, then applies SAGN/SLE and C\&S. LD prepares text features for graph aggregation; walk-text adds context from connected products. This complementary feature combination motivates the design, while the measured result belongs to the complete pipeline.\par

{\small \href{https://ogb.stanford.edu/docs/leader_nodeprop/#ogbn-products}{Source: official OGB node-property prediction leaderboard}.\par}

\section*{4. Parameter accounting}

\begin{longtable}{@{}p{4.88in}p{1.88in}@{}}

\toprule

Component & Parameters \\

\midrule

LD bert-base-uncased encoder & 109,482,240 \\

LD SAGN/SLE head and aggregation weight & 1,154,656 \\

EmbeddingGemma encoder & 307,581,696 \\

Final 1,636-input SAGN/SLE predictor & 5,010,014 \\

Correct \& Smooth & 0 \\

Total across listed components & \textbf{423,228,606} \\

\bottomrule\end{longtable}

{\small The total includes upstream feature-generation models and the final predictor; these parameters are not jointly trained. The LD rows sum to the published 110,636,896 parameters for LD+GIANT+SAGN+SCR. Training logs record the final SAGN count. [7]\par}

\newpage

\section*{5. Reproducibility and artifact checks}

The repository pins SAGN and LD revisions, includes the feature-integration patch, records external tensor hashes and sizes, and supplies the EmbeddingGemma reconstruction. Original generation-library versions, internal computation dtype, GPU, and resolved model revision were not captured; they remain unknown. [7]\par

\subsection*{5.1 Embedding reconstruction check}

A deterministic 128-node prefix was reconstructed using the pinned model and settings on an NVIDIA RTX PRO 4500 Blackwell. All \textbf{98,304 float32 values} matched the archived tensor rows exactly: mean cosine similarity 1.0, mean L2 difference 0.0, and maximum absolute difference 0.0. This verifies the tested prefix, not a full regeneration of every embedding. The comparison JSON and reconstruction environment are committed.\par

\subsection*{5.2 Evaluation replay from released predictions}

The v2-from-release script downloaded and verified all ten prediction tensors, recomputed all 27 validation candidates, recovered the identical configuration lock, and reproduced test accuracy \textbf{0.8749843545 ± 0.0004904255}. This is an evaluation replay from released predictions, not a fresh retraining of all ten models. Its compressed log and SHA-256 are preserved.\par

The public v2 release includes predictions, hashes, environment records, and logs. Small v1 provenance files are in the repository; large assets are mirrored in an archived release. Each seed log ends with an upstream SAGN NumPy-summary traceback after the final-stage checkpoint was saved. The wrapper verifies that checkpoint before continuing.\par

\section*{6. Conclusion}

The complete pipeline achieves \textbf{87.4984\% ± 0.0490\% test accuracy} across ten fresh seeds under the v2 validation-locked protocol. Released predictions and selection artifacts support auditing of the reported arithmetic. The result does not isolate a gain from walk-text or establish an official OGB rank. Any submission should declare external data and retain the full v1 historical disclosure.\par

\section*{References and supporting artifacts}

{\small [1] Hu et al. Open Graph Benchmark: Datasets for Machine Learning on Graphs. NeurIPS 2020. \href{https://ogb.stanford.edu/}{Source}.\par}

{\small [2] Sun et al. Scalable and Adaptive Graph Neural Networks with Self-Label-Enhanced Training. \href{https://github.com/skepsun/SAGN_with_SLE}{Source}.\par}

{\small [3] Shi et al. Label Deconvolution for Node Representation Learning on Large-scale Attributed Graphs against Learning Bias. \href{https://arxiv.org/abs/2309.14907}{Source}.\par}

{\small [4] Tan et al. WalkLM: A Uniform Language Model Fine-tuning Framework for Attributed Graph Embedding. NeurIPS 2023. \href{https://proceedings.neurips.cc/paper_files/paper/2023/hash/2ac879d1865475a7abc8dfc7a9c15c27-Abstract-Conference.html}{Source}.\par}

{\small [5] Huang et al. Combining Label Propagation and Simple Models Out-performs Graph Neural Networks. ICLR 2021. \href{https://openreview.net/forum?id=8E1-f3VhX1o}{Source}.\par}

{\small [6] Google. EmbeddingGemma documentation. \href{https://ai.google.dev/gemma/docs/embeddinggemma}{Source}.\par}

{\small [7] Experiment code, technical report, releases, and controlling tuning/protocol disclosure. \href{https://github.com/mshanthappa-hub/ld-walklm}{Source}.\par}

{\small \href{https://github.com/mshanthappa-hub/ld-walklm/blob/main/TUNING_AND_PROTOCOL.md}{Controlling protocol disclosure} · \href{https://github.com/mshanthappa-hub/ld-walklm/blob/main/report/technical_report.pdf}{Original technical report}\par}

\end{document}